\documentclass{report}
\usepackage{amsmath,amssymb}
\setlength{\parindent}{0mm}
\setlength{\parskip}{1em}
\begin{document}
\begin{center}
\rule{6in}{1pt} \
{\large Abdulrahman Manea \\
{\bf A Massively Parallel Semicoarsening Multigrid for 3D Reservoir Simulation on Multi-core and Multi-GPU Architectures}}

Green Earth Sciences Building \\ Room 065 \\ Stanford University \\ Stanford \\ CA 94305-2220
\\
{\tt amanea@stanford.edu}\\
Tchelepi Hamdi\end{center}

In this work, we have designed and implemented a massively parallel
version of the Semicoarsening Black Box Multigrid Solver [1], which is
capable of handling highly heterogeneous and anisotropic 3D reservoirs,
on a parallel architecture with multiple GPU�s. For comparison purposes,
the same algorithm was also implemented on a shared-memory multi-core
parallel architecture using OpenMP. The parallel implementation exploits
the parallelism in every module of the original Multigrid algorithm,
including both setup stage and solution stage, without modifying the
original algorithm basic steps. The benefits of this approach are
twofold: maintaining the inherent strong linear convergence of the serial
Multigrid algorithm, and making advantage of the shared-memory
architecture to minimize the need for communication.

The design of the algorithm uses a combination of plane relaxation and
semicoarsening to efficiently handle anisotropies in 3D, [2]. Sense the
z-direction in most reservoir models is a direction of strong-coupling
compared to the x- and y- directions, semicoarsening is employed in the
z-direction, and plane relaxation is used for relaxation on x-y planes.
Besides solving 2D-systems for plane-relaxation, during the setup stage,
a set of 2D systems must be also solved on each multigrid level to get an
approximate representation of the exact prolongation operator described
in [1]. For solving both types of 2D systems, we used a parallel version
of the 2D standard-coarsening operator-induced multigrid [3]. To be able
to handle problems involving high anisotropies in the x- and y-
directions, we used alternating line-relaxation with zebra ordering to
parallelize across multiple line solves. For the coarsest-solve, we use
multicolor Gauss-Seidel algorithm, with four colors to handle the
nine-point stencils of the coarsest structured grid.

There are several differences in the 2D-solver requirements between the
setup-stage systems and the solution-stage systems. The setup-stage
systems are independent and directly parallelizable. On the other hand,
for parallelizing the solution stage systems, i.e. plane-relaxation,
red-black coloring must be used. Since the convergence of the overall
semicoarsening multigrid algorithm is very sensitive to the quality of
the prolongation weights computed by the 2D-solves during the setup
stage�which are in turn are used to build the coarse-grid operators in
the standard Galerkin approach� the setup stage 2D-solves are carried out
to tight residual tolerances using several multigrid cycles. On the other
hand, for plane relaxation, the accuracy requirements are much looser,
where a single multigrid cycle is found to be good enough for smoothing
most of the high-frequency error modes.

To have coalesced global memory access pattern on the GPU, the Diagonal
Sparse-Matrix Format [4] was used to store the system matrices at all the
levels. In both the 3D Semicoarsening Multigrid and the 2D Multigrid
plane relaxation, V-cycling was used to avoid spending more time at
coarser levels and thus affecting the parallel efficiency. To minimize
the expensive communication through PCIe between the host and a GPU and
amongst GPU�s, every 2D solve is explicitly handled by a single GPU,
where we avoid splitting a single 2D solve between multiple GPU�s.

Due to the inherent granularity difference between the GPU threads and
the Open-MP threads, each GPU thread is typically assigned one
computational point, whereas every relatively coarse OpenMP thread is
typically assigned several computational points. In addition, this
inherent granularity difference also affects the design of the
tridiagonal solves during the line-relaxation. Since OpenMP threads are
coarse (as they are mapped to one physical core on the host), it is
suffices to assign a line or a small group of lines to a single thread,
where the tridiagonal systems can be efficiently handled using the serial
Thomas Algorithm. However, on the GPU, this approach would be very
inefficient, as it involves relatively large serial computational tasks,
and would not generate enough parallel work to utilize most of the
available hardware resources. Thus, the approach taken was to assign
every computational point to a thread, where threads operates in two
stages, namely, a setup stage and a solution stage. In the setup stage,
the threads work on preparing the tridiagonal systems of one red/black
zebra line sweep, in parallel. Then, in the solution stage, the threads
solve those tridiagonal systems in one batch in parallel using NVIDIA's
CUSPARSE LIBRARY [5], which in turn uses a combination of both the Cyclic
Reduction and Parallel Cyclic Reduction algorithms.

The two versions were tested using various highly heterogeneous problems
derived from SPE10 Second Dataset and varying in size from tens of
millions of cells to hundreds of millions of cells. For problems with
sizes large enough to saturate the GPU resources, the Multi-GPU
implementation is found to be faster than the OpenMP implementation
running on 12 Intel � Xeon � X5650 2.66GHz cores. In addition, the
inherent serial nature of multiplicative multigrid, along with the
approach taken to minimize the communication through PCIe, is found to
limit the scalability beyond 3-4 GPU�s.
\\
{\bf{References:}}
\\
\begin{enumerate}
\item S. Schaffer. A semicoarsening multigrid method for elliptic partial
differential equations with highly discontinuous and anisotropic
coefficients.SIAM J. Sci. Comput, 20(1):228�242, 1998.\\
\item Dendy, Jo E., et al. "Multigrid methods for three-dimensional
petroleum reservoir simulation." Tenth SPE Symposium on Reservoir
Simulation. 1989.\\
\item Alcouffe, R. E., et al. "The multi-grid method for the diffusion
equation with strongly discontinuous coefficients." SIAM Journal on
Scientific and Statistical Computing 2.4 (1981): 430-454.\\
\item N. Bell, M. Garland, "Efficient Sparse Matrix-Vector Multiplication
on CUDA", NVIDIA Technical Report, NVR-2008-004, (2008).\\
\item NVIDIA Corporation, "NVIDIA CUDA CUSPARSE Library", [Online].
Available: http://docs.nvidia.com/cuda/index.html
\end{enumerate}


\end{document}
